\chapter{Overview}
\label{ch:overview}

This library holds five synthesizable Verilog floating-point units in the
IEEE-754 binary format. Each \emph{unit} takes two operands with an input valid
strobe, and returns a result with an output valid strobe. Table~\ref{tab:units} lists the units and
their latencies in clock cycles (Section~\ref{sec:protocol}).

\begin{threetable}{Floating-Point Units}{tab:units}{Unit}{Latency}{Function}{1.6cm}{4.2cm}
\sig{fadd} & 5 & Addition (Chapter~\ref{ch:fadd}). \\ \hline
\sig{fmul} & 4 & Multiplication (Chapter~\ref{ch:fmul}). \\ \hline
\sig{fdiv} & $\mathtt{MAN\_W}+6$ (30 for binary32) & Division (Chapter~\ref{ch:fdiv}). \\ \hline
\sig{fmac} & 8, from the last product & Sum of a stream of products (Chapter~\ref{ch:fmac}). \\ \hline
\sig{fpu} & 5 for addition, subtraction and multiplication; as \sig{fdiv} for division & Addition, subtraction, multiplication and division (Chapter~\ref{ch:fpu}). \\
\end{threetable}

The units share the modules described in Chapter~\ref{ch:lib}. The remaining
chapters cover the make targets and variables (Chapter~\ref{ch:make}), simulation
and tests (Chapter~\ref{ch:sim}), FPGA compilation (Chapter~\ref{ch:fpga}), ASIC
synthesis (Chapter~\ref{ch:syn}) and the software model (Chapter~\ref{ch:sw}).

\section{Features}
\label{sec:features}

\begin{itemize}
\item Parameterizable word and exponent widths; the default is binary32.
\item Subnormal operands and results.
\item The five IEEE-754 rounding modes, selected per operation by \sig{rnd\_mode}
  (Table~\ref{tab:rnd_mode}).
\item Optional handling of NaN and infinity operands.
\item One operation per clock cycle in \sig{fadd} and \sig{fmul}, and in any mix
  of additions, subtractions and multiplications in \sig{fpu}.
\item Iterative division.
\item Accumulation of a stream of products in \sig{fmac}, with one product per
  clock cycle.
\item Simple interface with an input and an output valid strobe.
\item Asynchronous, active-high reset.
\end{itemize}

The units have no exception flag outputs: they do not report invalid operation,
division by zero, overflow, underflow or inexact results.

\section{Number Format}
\label{sec:format}

A word of \sig{DATA\_W} bits has a sign bit, an exponent field of \sig{EXP\_W}
bits and a fraction of $\mathtt{DATA\_W}-\mathtt{EXP\_W}-1$ bits, as
Figure~\ref{fig:format} shows. The \emph{mantissa} is the fraction with its
hidden bit on top, so it has $\mathtt{MAN\_W}=\mathtt{DATA\_W}-\mathtt{EXP\_W}$
bits. The hidden bit is 1 for a normal number and 0 for a subnormal number or
zero. Table~\ref{tab:classes} lists the encodings.

\begin{figure}[H]
\centering
\begin{tikzpicture}[x=1mm,y=1mm]
  \draw[fill=ug-head] (0,0) rectangle (10,8) node[midway,font=\small] {sign};
  \draw[fill=ug-block] (10,0) rectangle (50,8) node[midway,font=\small] {exponent};
  \draw[fill=white] (50,0) rectangle (130,8) node[midway,font=\small] {fraction};
  \node[font=\footnotesize,above] at (5,8) {1};
  \node[font=\footnotesize\ttfamily,above] at (30,8) {EXP\_W};
  \node[font=\footnotesize\ttfamily,above] at (90,8) {DATA\_W$-$EXP\_W$-$1};
\end{tikzpicture}
\caption{Word Format}
\label{fig:format}
\end{figure}

\begin{threetable}{Encodings}{tab:classes}{Class}{Exponent field}{Fraction}{3.2cm}{3.6cm}
Zero & All zeros & Zero \\ \hline
Subnormal & All zeros & Not zero \\ \hline
Normal & Neither all zeros nor all ones & Any \\ \hline
Infinity & All ones & Zero \\ \hline
NaN & All ones & Not zero \\
\end{threetable}

When a unit generates a NaN, the sign is 0, the exponent field is all ones, the
top fraction bit is 1 and the other fraction bits are 0. A NaN operand's sign and
fraction are not propagated.

\section{Common Parameters}
\label{sec:params}

All units have the parameters of Table~\ref{tab:common_params}. This guide writes
\sig{MAN\_W} for $\mathtt{DATA\_W}-\mathtt{EXP\_W}$, the mantissa width including the
hidden bit; it is a local parameter of each unit. The test suite runs with
\sig{DATA\_W}~=~32 and \sig{EXP\_W}~=~8 (Section~\ref{sec:testfloat}).

\begin{paramtable}{Common Parameters}{tab:common_params}
DATA\_W & 32 & Word width in bits. \\ \hline
EXP\_W & 8 & Exponent field width in bits. \\ \hline
SPECIAL\_CASES & 1 & 1: handle NaN and infinity operands (and, in division, a zero divisor) as described in each unit's chapter. 0: remove this handling, so these operands go through the datapath as ordinary numbers and the result is not defined by this guide. \\
\end{paramtable}

\begin{twotable}{Rounding Modes}{tab:rnd_mode}{\sig{rnd\_mode}}{Rounding}{2.6cm}
\sig{0} & To nearest, ties to even. The IEEE-754 default. \\ \hline
\sig{1} & Toward zero. \\ \hline
\sig{2} & Toward $-\infty$. \\ \hline
\sig{3} & Toward $+\infty$. \\ \hline
\sig{4} & To nearest, ties away from zero. \\ \hline
\sig{5}--\sig{7} & As \sig{0}. \\
\end{twotable}

The encoding is that of the RISC-V \sig{frm} field. On overflow, a mode that rounds
toward zero for the sign of the result gives the largest finite number instead of
infinity. An exact zero from operands of opposite signs is $+0$, or $-0$ in mode 2.

\section{Valid Protocol}
\label{sec:protocol}

All units use the same signals to start an operation and return its result:

\begin{enumerate}
\item After reset, drive the operands and raise \sig{valid\_i} for one clock cycle.
  The unit samples \sig{valid\_i} on the rising edge of \sig{clk}.
\item Wait for \sig{valid\_o}, which is high for one clock cycle per result.
\item Read \sig{res} in the clock cycle in which \sig{valid\_o} is high. \sig{res} is
  valid only in that cycle.
\end{enumerate}

The \emph{latency} $L$ is the distance in clock cycles between the cycle in which
\sig{valid\_i} is high (cycle 0) and the cycle in which \sig{valid\_o} is high (cycle
$L$). All latencies in this guide were measured in simulation, with Icarus
Verilog, and did not change between normal, subnormal, zero, infinite and NaN
operands.
Figure~\ref{fig:timing_pipe} shows an operation with $L=5$.

\begin{figure}[H]
\centering
\begin{tikztimingtable}[timing/.style={x=3.2mm,y=3mm},timing/d/text/.style={font=\scriptsize}]
cycle & 2D{0} 2D{1} 2D{2} 2D{3} 2D{4} 2D{5} 2D{6} \\
\sig{clk} & 7{HL} \\
\sig{valid\_i} & 2H 12L \\
\sig{op\_a}, \sig{op\_b} & 2D{op} 12U \\
\sig{valid\_o} & 10L 2H 2L \\
\sig{res} & 10U 2D{res} 2U \\
\end{tikztimingtable}
\caption{Timing Diagram, Operation With a Latency of 5}
\label{fig:timing_pipe}
\end{figure}

All units sample the operands in the \sig{valid\_i} cycle. The unit chapters
state what differs: when the next operation can start.

\section{Source Files}
\label{sec:files}

Each unit is in its own directory, with its sources in \file{<unit>/hw/src}. The
shared modules are in \file{lib/src}. Table~\ref{tab:files} lists the Verilog
files needed to compile each unit, as in its \file{hw/hw.mk}.

\begin{twotable}{Source Files of Each Unit}{tab:files}{Unit}{Files}{1.6cm}
\sig{fadd} & \file{fadd/hw/src/fadd.v}, \file{fadd/hw/src/fadd\_core.v}, and the \file{lib/src} files \file{unpack.v}, \file{cls.v}, \file{norm.v}, \file{round.v}, \file{pack.v} \\ \hline
\sig{fmul} & \file{fmul/hw/src/fmul.v}, \file{fmul/hw/src/fmul\_core.v}, and the same \file{lib/src} files \\ \hline
\sig{fdiv} & \file{fdiv/hw/src/fdiv.v}, \file{fdiv/hw/src/fdiv\_core.v}, and the same \file{lib/src} files \\ \hline
\sig{fmac} & \file{fmac/hw/src/fmac.v}, \file{fmac/hw/src/fmac\_core.v}, and the \file{lib/src} files except \file{norm.v} \\ \hline
\sig{fpu} & \file{fpu/hw/src/fpu.v}, \file{fadd/hw/src/fadd\_core.v}, \file{fmul/hw/src/fmul\_core.v}, \file{fdiv/hw/src/fdiv\_core.v}, and the same \file{lib/src} files as \sig{fadd} \\
\end{twotable}

Each unit directory also has \file{hw/sim} (Chapter~\ref{ch:sim}), \file{hw/fpga}
(Chapter~\ref{ch:fpga}) and \file{hw/syn} (Chapter~\ref{ch:syn}).

\section{Common Interface Signals}
\label{sec:signals}

Table~\ref{tab:common_sigs} lists the signals of all the units. \sig{fpu} adds
\sig{funct} (Table~\ref{tab:fpu_sigs}). All outputs are registered.

\begin{sigtable}{Common Interface Signals}{tab:common_sigs}
clk & Input & 1 & Clock. The unit uses its rising edge. \\ \hline
rst & Input & 1 & Asynchronous reset, active high. \\ \hline
valid\_i & Input & 1 & Input valid: the operands of a new operation are on the inputs. \\ \hline
valid\_o & Output & 1 & Output valid: the result is valid in this clock cycle. \\ \hline
rnd\_mode & Input & 3 & Rounding mode (Table~\ref{tab:rnd_mode}), sampled with the operands. \\ \hline
op\_a & Input & DATA\_W & First operand. \\ \hline
op\_b & Input & DATA\_W & Second operand. \\ \hline
res & Output & DATA\_W & Result. \\
\end{sigtable}
